\section{Problem Description}\label{sec:problem}

A manufacturer must complete a set of jobs with identical internal parallel
machines and an external production provider. Before operations begin, the
manufacturer decides whether each job is processed internally or outsourced.
Internally retained jobs are assigned, sequenced, and timed on the available
machines. An outsourced job no longer consumes internal processing or setup time;
its production and delivery consequences are assumed to be included in the
provider's quoted baseline cost. Each job is processed exactly once by one of
these two modes.

Table~\ref{tab:core-notation} summarizes the notation used throughout the model.
\begin{table}[t]
    \centering
    \caption{Core problem notation.}
    \label{tab:core-notation}
    \begin{tabularx}{\textwidth}{@{}lX@{}}
        \toprule
        Symbol & Definition \\
        \midrule
        $\mathcal J,\mathcal M$ & Sets of jobs and identical internal machines \\
        $p_j,[d_j^e,d_j^l]$ & Processing time and preferred completion window of job $j$ \\
        $s_{ij},\kappa_{ij}$ & Setup time and setup cost when job $j$ follows $i$ \\
        $\phi_j(t)$ & Internal completion-time penalty of job $j$ at time $t$ \\
        $b_j$ & Baseline cost quoted for outsourcing job $j$ \\
        $B(O)$ & Total baseline outsourcing cost of job set $O$ \\
        $G(B(O))$ & Actual payment after applying the aggregate outsourcing tariff \\
        $\mathcal R$ & Set of feasible elementary single-machine schedules \\
        $c_r^{\mathrm{in}}$ & Minimum setup and timing cost of internal schedule $r$ \\
        \bottomrule
    \end{tabularx}
\end{table}

\subsection{Internal production}

Let $\mathcal J=\{1,\ldots,n\}$ be the set of jobs and
$\mathcal M=\{1,\ldots,m\}$ the set of identical internal machines. All jobs and
machines are available at time zero. Internal processing is nonpreemptive, and a
machine processes at most one job at a time.

Job $j\in\mathcal J$ has processing time $p_j$ and due window
$[d_j^e,d_j^l]$. If job $j$ follows job $i$ on the same machine, the transition
requires setup time $s_{ij}$ and incurs setup cost $\kappa_{ij}$. Index $0$
denotes the initial state of a machine; hence $s_{0j}$ and $\kappa_{0j}$ are the
initial setup time and cost of job $j$. Idle time is allowed between consecutive
jobs. The setup times satisfy the directed triangle inequality
\begin{equation}
    s_{ik}\le s_{ij}+s_{jk},
    \qquad
    i\in\mathcal J\cup\{0\},\quad j,k\in\mathcal J,
    \label{eq:setup-time-triangle}
\end{equation}
for pairwise distinct indices. No corresponding triangle inequality is imposed on
the setup costs.

The completion-time penalty of job $j$ is a nonnegative, finitely represented
piecewise-linear function $\phi_j(t)$. The principal case is weighted
earliness--tardiness with a due window,
\begin{equation}
    \phi_j(t)
    =
    \alpha_j[d_j^e-t]^+
    +
    \beta_j[t-d_j^l]^+,
    \qquad [x]^+:=\max\{x,0\}.
    \label{eq:due-window-penalty}
\end{equation}
The pricing algorithms require sums, translations, lower envelopes, and interval
minima of these functions to be evaluated exactly. The representation therefore
also includes due dates as zero-width windows and other finitely represented
piecewise-linear early and late completion costs.

An internal machine schedule is an elementary sequence
$r=(j_1,\ldots,j_q)$ of distinct jobs. Let $j_0=0$ and $C_{j_0}=0$. The minimum
cost associated with this order is
\begin{equation}
    c_r^{\mathrm{in}}
    =
    \sum_{h=1}^{q}\kappa_{j_{h-1}j_h}
    +
    \min_{\boldsymbol C}
    \left\{
        \sum_{h=1}^{q}\phi_{j_h}(C_{j_h})
        :
        C_{j_h}\ge
        C_{j_{h-1}}+s_{j_{h-1}j_h}+p_{j_h}
    \right\}.
    \label{eq:internal-column-cost}
\end{equation}
The inequalities in \eqref{eq:internal-column-cost} permit intentional idle time.
Thus, an internal column specifies a job order, while its coefficient already
contains the optimal completion times for that order.

\subsection{Aggregate outsourcing tariff}

Job $j$ has baseline outsourcing cost $b_j\ge0$. For an outsourced set
$O\subseteq\mathcal J$, its total baseline outsourcing cost is
\begin{equation}
    B(O)=\sum_{j\in O}b_j.
    \label{eq:baseline-outsourcing-cost}
\end{equation}
The actual payment is $G(B(O))$, where $G:\mathbb R_+\rightarrow\mathbb R_+$ is
continuous, nondecreasing, and concave, with $G(0)=0$. The baseline model uses a
finite piecewise-linear tariff. Let $\mathcal L=\{1,\ldots,L\}$ be its segments.
Segment $\ell$ applies on
$[\underline q_\ell,\overline q_\ell]$ and has affine expression
\begin{equation}
    G(q)=\gamma_\ell q+\delta_\ell
    \qquad
    q\in[\underline q_\ell,\overline q_\ell].
    \label{eq:outsourcing-tariff-segment}
\end{equation}
Concavity implies nonincreasing marginal rates $\gamma_\ell$, while
nondecreasingness requires $\gamma_\ell\ge0$. The no-discount case is the
single-segment tariff $G(q)=q$.

The tariff is aggregate rather than job-separable. Consequently, the incremental
payment for outsourcing job $j$ depends on the baseline cost already assigned to
the provider. Concavity captures quantity discounts in the studied problem, but
Section~\ref{subsec:sp1} explains why the outsourcing-column formulation can also
handle a broader class of outsourcing cost functions.

\subsection{Integrated optimization problem}

A feasible solution consists of at most $m$ mutually disjoint elementary internal
schedules and one outsourced set. The outsourced set is disjoint from all internal
schedules, and every job appears exactly once, either in one internal schedule or
in the outsourced set. The objective is
\begin{equation}
    \min
    \left\{
        \sum_{r\in\mathcal R^\star}c_r^{\mathrm{in}}
        +
        G(B(O))
    \right\},
    \label{eq:original-problem}
\end{equation}
where $\mathcal R^\star$ is the selected collection of internal schedules and
$O$ is the outsourced job set. The formulations below use covering rows, matching
the restricted master problems implemented in the exact algorithm.
